\scanpage{107}
Soit $J$ un ensemble ordonné, on s'intéresse aux plongements d'ensembles ordonnés
\[
(150)\qquad J\mathrel{\overset{\varphi}{\hookrightarrow}}\Delta_x\unsure{\xi}\simeq\mathcal P(I(x))^\circ
\]
Un tel plongement donne une ``fonction codimension'' sur $J$, en composant, \inserted{soit $d$,} et \deleted{dire que $\varphi$ est un plongement d'ens. ordonnés signifie que}\inserted{on suppose donné déjà $d$ d'avance sur $J$} satisfaisant
\[
(151)\qquad\text{Si }x<y\quad\text{on a}\quad\unsure{d|c}(x)<c(y),
\]
et on se borne aux $\varphi$ qui sont des \unsure{plongts|plongements} ordonnés, compatibles avec \unsure{foncts} codim sur $J$ et sur $\Delta_x$.
\deleted{On suppose de plus que $J$ a un plus petit élément (ce qui correspond donc à une codim. maximale \deleted{qui est} soit $d(J)$) — quitte à le rajouter, s'il n'y est déjà!}

Considérons, \inserted{pour tout $\unsure{q|d}\in\mathbb N^*$,} la multiplicité analytique formée des couples $(x,\unsure{V})$, où $x$ est un pt de $D$ et $V$ un germe de sous-multiplicité de $X$ attaché à $D$, de codim $\unsure{q|d}$. On \deleted{voit} voit aisément que ce sont \inserted{bien} là les points d'une multiplicité analytique que je note $D_{\unsure{q|d}}$, et on a un morphisme
\margin{106}
